% Folder 144, batch 01 (archivists' pages 1-20).
% Pass: Opus 5.5 (claude-opus-5-5), 2026-10-03 — first pass, unchecked against the pages by a human.
%
% Page 1 is a cover sheet in his hand (its words are the section heading);
% page 11 is blank; neither gets a page marker. Pages 9-10 are a letter from
% Barry Mazur (typed, with his green-ink additions, and a handwritten PS),
% third-party correspondence: summarised in notes, not transcribed (RIGHTS.md).
\documentclass[11pt,a4paper]{article}
\input{../preamble/grothendieck}

\folder{144}
\batch{1}
\pages{1}{20}
\foldertitle{[Suite autour de Teichmüller dont] Teichmülleries, 1981-1982 : notes manuscrites (1981-1983, s.d.), lettres (1981, s.d.).}
\dating{1981-1983}
\watermark{Édition de démonstration}

\begin{document}

\section{Discussion avec Deligne + lettre Barry Mazur}
\note{Titre de sa main sur le feuillet de garde (p.~1), qui ne porte rien
d'autre.}

\page{2}
Conversations avec Deligne du 30.5 et 5.6.

1) Pour groupes finis opérant sur surfaces~: utiliser «~th.\ de Smith~»
(suites spectrales) pour prouver que le lieu des points fixes est une
variété homologique du pt de vue local -- donc \ill{} loc.\ connexe,
donc connexe par arcs…

2) Exemple de Bing d'une involution de $B^3$ ayant un cercle de pts fixes
sauvage (sphère \add{\uncertain{genre}} \uncertain{cornue} d'Alexander)
mais, dans la \ill{}~: $S^2=\partial B^3$ \ill{} la conj.\ complexe.

3) Il doit y avoir des actions \add{simpliciales} de groupes finis sur
$B^n$ ($n\geq{}$\struck{\ill{}} $5$~??) qui ne sont pas le cône d'une
action sur $S^{n-1}=\partial B^n$. Lié au fait qu'une
\struck{\ill{}} action (d'un groupe fini sur une sphère avec pt fixe)
\struck{\ill{}} n'est pas nécess.~! contenue dans un conjugué
\add{topologique} du groupe orthogonal. (J'ai pas bien compris
comment….)

4) Nécessité de faire intervenir le \uncertain{fermé universel}
(d'ailleurs dans $\bar{T}$) pour décrire \struck{les gps} \add{la}
structure «~arithmétique~» du groupe fondamental d'une courbe lisse
propre connexe sur \uncertain{$k$ alg.\ clos}. Question de définir
cohomologiquement les intersections de «~chemins~» (ou sans théorie des
revêtements)

\drawing{Deux chemins orientés entre points marqués, l'un coupant
l'autre.}

Données théories cohomologiques d'une courbe dans le \ill{} ……

\drawing{Deux autres figures de chemins orientés à points marqués, avec
une boucle, séparées par «~ou~».}

\page{3}
5) Compactification des topos modulaires $M_{g,\nu}$ et structure à
l'infini. Il semble qu'on ait $\pi_1(\widehat{M_{g,\nu}})=1$ si $g\geq 2$
(si $g=1$, $\nu=0$, le $\pi_1$ doit être cyclique \uncertain{fini}
d'ordre fini divisant~12…), et $\pi_1(M^{\infty}_{g,\nu})=1$ itm.

\marginal{*} Déf.\ de $\widehat{M}_{g,\nu}$~: schémas $X/S$ munis de
sections $(f_1,\dots,f_\nu)$, $X/S$ \uncertain{schémas propres}
(\uncertain{pas plats}~?) et plats sur $S$, de prés.\ finie, à fibres
$X_s$ \uncertain{géom.} connexes \uncertain{réduites}, n'ayant comme
singularités que des pts doubles ordinaires, et les $f_{is}$ étant des
pts lisses, enfin le groupe des automorphismes de
$(X_s,f_{1s},\dots,f_{\nu s})$ étant fini (ce qui signifie que
\struck{les} chaque composante irréductible de genre~0 de $X_s$ a au
moins trois points marqués -- soit par intersection avec les autres
comp.\ irréd., soit comme pts $f_{is}$…)

\marginal{*} Question~: si $\nu=0$, $X_s$ \add{lisse} est-il contenu
dans sa jacobienne, extension d'une VA par un tore~?

La vérification des critères valuatifs \marginal{*} semble
\uncertain{étrange}.

Les générateurs «~twist~» de $\mathcal{T}_{g,\nu}$ sont les lacets
autour des comp.\ irréductibles de $M^{\infty}_{g,\nu}$….

\page{4}
6) Nécessité, en dim${}>1$, d'une description plus délicate de la
«~structure à l'infini~» d'un schéma lisse connexe
\add{\uncertain{séparé}} sur $k$ alg.\ clos.

$X$

\begin{tikzcd}[column sep=small, row sep=small, nodes={font=\scriptsize}]
1 \arrow[r] & \pi_1^{\infty}(X) \arrow[r, "i"] & \pi_1(X) \arrow[r] & \pi_1(\hat X) \arrow[r] & 1 \\
1 \arrow[r] & \pi_1^{\infty}(\mathcal{U}) \arrow[r] \arrow[u, "\alpha"] & \pi_1(\mathcal{U}) \arrow[r] \arrow[u] & \pi_1(\mathcal{X}) \arrow[r] \arrow[u] & 1 \\
 & & & \pi_1(Y) \arrow[u, no head, "\simeq"'] &
\end{tikzcd}

\note{Au-dessus de $\pi_1(\hat X)$, la page ajoute «~compactifié~».
Devant la seconde ligne, \struck{$\hat X\supset Y$}.}

$\mathcal{U}\subset\mathcal{X}$ \add{complété formel de $\hat X$ le long
de $Y$}, au-dessus de $X\subset\hat X\supset Y$
($=\hat X\setminus X$).

$\pi_1^{\infty}(X)$ est le sous-groupe \add{fermé} invariant de
$\pi_1(X)$ engendré par $\alpha(\pi_1^{\infty}(\mathcal{U}))$, mais
$\alpha$ n'est pas nécess.\ surjectif. Il faudrait expliciter,
\add{\uncertain{comme on a dit}}, en termes de «~combinatoire~»
$T=\varprojlim_n$ \ill{} et de la structure de $Y$ (qu'on pourrait
supposer \add{\ill{}} diviseur à croisements normaux) le groupe
$\pi_1^{\infty}(\mathcal{U})$ et l'action de Galois dessus, et
\add{l'hom.\ ext.} exiger que ($i\circ\alpha:
\pi_1^{\infty}(\mathcal{U})\to\pi_1(X)$) commute à l'action de Galois…

On doit pouvoir \add{par récurrence} faire une réduction au cas
$\dim X=2$ ….

7) Sur l'espace de Teichmüller \add{complexe}
$\widetilde{M}^{\mathrm{an}}_{g,\nu}$ ($=\mathcal{M}^{\mathrm{an}}_{g,\nu}$)
(rev.\ universel de \struck{la multiplicité} \add{\ill{}}
$M^{\mathrm{an}}_{g,\nu}(\mathbb{C})$)

\page{5}
\note{Suite du point 7) de la p.~4, sur l'espace de Teichmüller.}
on \uncertain{n'a pas défini jusqu'à présent} de métrique
\uncertain{riemannienne} \uncertain{intrinsèque}, mais bien des
métriques \uncertain{finsleriennes}. On espère trouver une métrique
riemannienne \ill{} \emph{intégration}, qui permettrait par un argument
géométrique simple (cf.\ l'existence \add{\uncertain{et unicité}} de la
géodésique qui joint deux pts) de prouver que tout groupe fini
d'automorph.\ complexes de $\mathcal{M}_{g,\nu}$ admet un pt fixe.
Deligne pense que \add{c'est vrai que} tout groupe fini de
$\mathcal{T}_{g,\nu}$ \uncertain{admet} un pt fixe dans $M_{g,\nu}$, et
il ne doit pas être \ill{} \ill{}.

\drawing{Points $x_0, x_1, gx_0, gx_1, g^2x_0, g^2x_1, \dots,
g^{N-1}x_0, g^{N-1}x_1$ et $x_2, gx_2, \dots, g^{N-1}x_2$, disposés en
polygones emboîtés qui se resserrent~; une flèche désigne le point
central, légendé «~pt fixe par \ill{}~».}

8) Mumford connaîtrait la présentation des $\mathcal{T}_{g,\nu}$ via les
générateurs habituels).

L.~Bers (Columbia \uncertain{Sc.}) ferait de la «~chirurgie~» de
surfaces de Riemann pour \struck{déterminer} établir la compactification
des $\mathcal{M}_{g,\nu}$.

\page{6}
9) Le cône sur une variété compacte $S$ peut être une variété~: bord,
sans que $S$ soit une sphère.

Exemple~: considérons le $\uncertain{4}$-groupe $\Gamma\subset S^3$
d'ordre~120, \struck{\ill{}} et
$\mathfrak{G}^{+}_{5}\subset S^3/\pm1$ ($=SO(3,\mathbb{R})$), image de
$\Gamma$, $\mathfrak{G}^{+}_{5}=\mathfrak{A}$. Considérons
$S^3/\Gamma\simeq SO(3,\mathbb{R})/\mathfrak{A}\simeq{}$classification
des icosaèdres \add{\uncertain{réguliers}} inscrits dans une sphère
donnée. Ce n'est pas une sphère, \struck{\ill{}} sa suspension, qui a
deux pts singuliers, n'est pas une variété (elle a \uncertain{en} ces
pts, le groupe fondamental local est $\simeq\Gamma$, donc non trivial)
Mais \struck{elle a un} \add{\ill{}} sa \uncertain{double suspension}
\add{\uncertain{(suspension itérée)}} \struck{est \ill{}} est
\struck{la suspension} le produit par une droite et
\struck{\ill{}} est une variété (?), et $\simeq S^5$, et $\simeq$
\ill{} \add{\uncertain{$S^5$}}.

\marginal{Edwards, Exposé Bourbaki 1977--78 (Cannon, 515)~: $X$ variété
de dim.\ $n$, $\mathbb{Z}$-sphère homologique, sa double suspension est
une sphère.}

\uncertain{En fait}, sa double suspension est une variété -- y compris
aux deux «~pôles~» ….

Pour donner \add{\uncertain{les limites de nature \ill{} topolog.}} des
conditions \struck{\ill{}} sur un espace \uncertain{à
stratification} \struck{\ill{}} \add{\uncertain{conique}}
(\uncertain{choisissant} \struck{l'interprétation} \add{\ill{}} qui
\uncertain{l'implique} \add{\uncertain{du fait que}} certaines
\uncertain{espaces}, suivant des sphères, \struck{\ill{}}, il
faudrait d'après Deligne se placer dans un cadre

\page{7}
tel que le PL (ou d'autres milieux, \ill{} p.\,ex.), le cadre topologique
ne s'y prêtant pas. Il semble néanmoins que dans le cas des variétés
loc.\ finies \emph{régulières}, les difficultés disparaissent,
\struck{grâce au fait que}

\marginal{à vérifier avec soin …}

10) Les groupes apparaissant comme groupes d'automorphismes de «~cartes~»
régulières \struck{\ill{}} \emph{connexes} (loc.\ finies~?),
\struck{\ill{}} \add{variétés} \uncertain{en fait}, \uncertain{sont tous}
des groupes de Coxeter. \marginal{id} Ceux parmi eux qui sont finis sont
bien connus (cf.\ livre Bourbaki) -- il semble qu'ils correspondent tous~:
des sphères -- avec la décomposition «~simpliciale~» évidente, liée à la
géométrie des groupes de Coxeter. D'après Deligne, toutes les «~cartes~»
régulières \add{simplement connexes} (finies ou non) (\uncertain{données}

\page{8}
figurent dans Bourbaki (pour groupes \uncertain{coniques} finis etc.) --
tous au moins dans le cas \uncertain{non singulier}. Dans le cas infini,
(le cas semblerait lié \uncertain{à l'hyperbolique}~?)~: celui où la
forme quadratique invariante est (de \uncertain{signe} \ill{}) de
signature $(1,\dots,1)$ -- ce qui signifie que la sphère unité se
décompose en deux composantes \emph{connexes simplement connexes}~:
c'est une des composantes qui \add{doit} jouer le rôle de la
\struck{\ill{}} variété, en prenant les traces dessus des cônes
simpliciaux (\ill{} de $\mathfrak{A}$) de l'un d'eux…

11) Lien entre l'icosaèdre et $E_8$~: la forme quadratique
\uncertain{de degré}~4 sur $\mathbb{Z}[\alpha]=\mathbb{Z}[T]/(T^2+T-1)$
$=$ anneau des entiers de
$\mathbb{Q}\bigl(\tfrac12(-1\pm\sqrt5)\bigr)=\mathbb{Q}(\sqrt5)$,
\marginal{?} composée avec la trace, donne la forme quadratique de
$E_8$ …. d'où syst.\ de racines de $E_8$…. Mais attention~: la
ramification \uncertain{en~5} …

\page{9}
\note{Lettre dactylographiée de Barry Mazur à Grothendieck, en anglais, sans date, avec des ajouts de sa main à l'encre verte. Elle répond aux questions d'une lettre de Grothendieck : un contre-exemple de Deligne à la question (a) (une extension non scindée d'un groupe fini par le revêtement universel de $\mathrm{SL}_2(\mathbb{R})$), la question (b), une involution de la boule $B^3$ à lieu fixe sauvage, les sous-groupes finis du groupe de Teichmüller ; et elle fixe sa visite (fin mai, début juin). Correspondance d'un tiers, non transcrite (RIGHTS.md).}

\page{10}
\note{Post-scriptum manuscrit de Mazur, au verso de la même feuille : ce qu'en disent Siebenmann et Kravitz (sous-groupes finis du groupe de Teichmüller et courbes algébriques ; automorphismes finis, théorie de Smith). Correspondance d'un tiers, non transcrite (RIGHTS.md).}

\page{12}
\marginal{1982}
\note{Le groupe noté ici $\mathfrak{S}$ est écrit d'une lettre gothique
cursive, la même qu'en $\mathfrak{S}_3$ p.~13~; pour le «~$\tau$~» cursif du groupe de
Teichmüller, $\mathcal{T}$. Les traits de ce premier diagramme sont sans
flèches sur la page.}

(1)

\begin{tikzcd}[column sep=small, row sep=small, nodes={font=\scriptsize}]
\mathcal{X}_{g,\nu\,\mathbb{Q}} \simeq M_{g,\nu+1,\nu\,\mathbb{Q}} \arrow[r, no head] \arrow[d, no head] & \mathcal{X}_{g,\nu\,\mathbb{C}} \simeq M_{g,\nu+1,\nu\,\mathbb{C}} \arrow[r, no head] \arrow[d, no head] & \mathcal{X}^{\mathrm{an}}_{g,\nu} \simeq M^{\mathrm{an}}_{g,\nu+1,\nu} \arrow[d, no head] \\
M_{g,\nu\,\mathbb{Q}} \arrow[r, no head] \arrow[d, no head] & M_{g,\nu\,\mathbb{C}} \arrow[r, no head] \arrow[d, no head] & M^{\mathrm{an}}_{g,\nu} \\
\mathbb{Q} \arrow[r, no head] & \mathbb{C} &
\end{tikzcd}

$X_{g,\nu}$ courbe alg./$\mathbb{C}$ type $g,\nu$, \uncertain{corr.}~:
$\xi_0\in\operatorname{Ob}\mathrm{Pt}(M_{g,\nu\,\mathbb{Q}};\mathbb{C})$
${}\simeq\operatorname{Ob}\mathrm{Pt}(M_{g,\nu\,\mathbb{C}}/\mathbb{C})$
${}\simeq\operatorname{Ob}\mathrm{Pt}(M^{\mathrm{an}}_{g,\nu})$,
$X_{g,\nu}\simeq(\mathcal{X}_{g,\nu})_{\xi_0}$~;
$x_0\in X_{g,\nu}(\mathbb{C})\simeq X^{\mathrm{an}}_{g,\nu}$,
$(\xi_0,x_0)$ \uncertain{corr.}~: \struck{$\xi_0$}
$\in\operatorname{Ob}\mathrm{Pt}(\mathcal{X}_{g,\nu\,\mathbb{Q}};\mathbb{C})$
${}\simeq\operatorname{Ob}\mathrm{Pt}(\mathcal{X}_{g,\nu\,\mathbb{C}}/\mathbb{C})$
${}\simeq\operatorname{Ob}\mathrm{Pt}(\mathcal{X}^{\mathrm{an}}_{g,\nu})$
au-dessus de $\xi_0$.

\textbf{Déf.} (1)
\[
\begin{cases}
\mathcal{T}^{+}_{g,\nu} = \pi_1(M^{\mathrm{an}}_{g,\nu},\xi_0)\\
\mathfrak{S}^{+}_{g,\nu} = \pi_1(\mathcal{X}^{\mathrm{an}}_{g,\nu},x_0)
  \simeq \mathcal{T}_{g,\nu+1,\nu} \subset \mathcal{T}_{g,\nu+1}\\
\pi_{g,\nu} = \pi_1(X_{g,\nu},x_0)\\
\mathcal{N}_{g,\nu} = \pi_1(M_{g,\nu\,\mathbb{Q}},\xi_0)\\
\mathcal{M}_{g,\nu} = \pi_1(\mathcal{X}_{g,\nu\,\mathbb{Q}},x_0)
\end{cases}
\]

d'où (2)

\begin{tikzcd}[column sep=small, row sep=small, nodes={font=\scriptsize}]
1 \arrow[r] & \pi_{g,\nu} \arrow[r] \arrow[d, no head, "="] & \mathfrak{S}^{+}_{g,\nu} \arrow[r] \arrow[d, "i"] & \mathcal{T}^{+}_{g,\nu} \arrow[r] \arrow[d, "i"] & 1 \\
1 \arrow[r] & \pi_{g,\nu} \arrow[r] & \mathrm{Aut}_{\mathrm{loc}}(\pi_{g,\nu}) \arrow[r] & \mathrm{Autext}_{\mathrm{loc}}(\pi_{g,\nu}) &
\end{tikzcd}

\struck{et} et, en prenant \uncertain{les complétés}~: (3)

\begin{tikzcd}[column sep=small, row sep=small, nodes={font=\scriptsize}]
1 \arrow[r] & \hat\pi_{g,\nu} \arrow[r] \arrow[d, no head, "="] & \hat{\mathfrak{S}}^{+}_{g,\nu} \arrow[r] \arrow[d, "?"] & \hat{\mathcal{T}}^{+}_{g,\nu} \arrow[r] \arrow[d, "?"] & 1 \\
1 \arrow[r] & \hat\pi_{g,\nu} \arrow[r] & \mathrm{Aut}_{\mathrm{loc}}(\hat\pi_{g,\nu}) \arrow[r] & \mathrm{Autext}_{\mathrm{loc}}(\hat\pi_{g,\nu}) &
\end{tikzcd}

\note{Les flèches verticales de (2) portent un petit signe, lu ici $i$
(«~injectif~» ?).}

\struck{\ill{}}

(4)
\[
\begin{cases}
\hat{\mathcal{T}}^{+}_{g,\nu} \overset{\mathrm{can}}{\simeq}
  \pi_1(M_{g,\nu\,\mathbb{C}},\xi_0) \simeq \pi_1(M_{g,\nu\,\bar{\mathbb{Q}}},\xi_0)\\
\hat{\mathfrak{S}}^{+}_{g,\nu} \overset{\mathrm{can}}{\simeq}
  \pi_1(\mathcal{X}_{g,\nu\,\mathbb{C}},x_0) \simeq \pi_1(\mathcal{X}_{g,\nu\,\bar{\mathbb{Q}}},x_0)
  \simeq \hat{\mathcal{T}}_{g,\nu+1,\nu} \subset \hat{\mathcal{T}}_{g,\nu+1}\\
\hat\pi_{g,\nu} \overset{\mathrm{can}}{\simeq}
  \pi_1(X_{g,\nu\,\mathbb{C}},x_0) \simeq \pi_1(X_{g,\nu\,\bar{\mathbb{Q}}},x_0)
\end{cases}
\]

(5)

\begin{tikzcd}[column sep=small, row sep=small, nodes={font=\scriptsize}]
 & 1 & 1 & & \\
1 \arrow[r] & \hat{\mathcal{T}}^{+}_{g,\nu} \arrow[r] \arrow[u] & \mathcal{N}_{g,\nu} \arrow[r] \arrow[u] & \Gamma_{\bar{\mathbb{Q}}/\mathbb{Q}} \arrow[r] & 1 \\
1 \arrow[r] & \hat{\mathfrak{S}}^{+}_{g,\nu} \arrow[r] \arrow[u] & \mathcal{M}_{g,\nu} \arrow[r] \arrow[u] & \Gamma_{\bar{\mathbb{Q}}/\mathbb{Q}} \arrow[r] \arrow[u, no head, "="'] & 1 \\
 & \hat\pi_{g,\nu} \arrow[r, no head, "="] \arrow[u] & \hat\pi_{g,\nu} \arrow[u] & & \\
 & 1 \arrow[u] & 1 \arrow[u] & &
\end{tikzcd}

\note{Une accolade sous les deux colonnes de gauche renvoie à (3).}

\struck{(cf.\ \ill{} $\hat\pi_{g,\nu}\to\mathcal{M}_{g,\nu}\to\mathcal{N}_{g,\nu}\to 1$)}

\textbf{NB} $\hat{\mathcal{T}}_{g,\nu}$, $\hat{\mathfrak{S}}_{g,\nu}$
s'obtiennent comme images inverses de
$\{1,\tau\}\subset\Gamma_{\bar{\mathbb{Q}}/\mathbb{Q}}$
($\tau\in\Gamma_{\bar{\mathbb{Q}}/\mathbb{Q}}$, conj.\ complexe) dans
$\mathcal{N}_{g,\nu}$, $\mathcal{M}_{g,\nu}$.

\textbf{NB} On a (via $\mathcal{N}_{g,\nu}\to\mathrm{Autext}_{\mathrm{loc}}(\hat\pi_{g,\nu})$
p.\,ex.)~: (7)

\begin{tikzcd}[column sep=small, row sep=small, nodes={font=\scriptsize}]
\mathcal{T}^{+}_{g,\nu} \arrow[r] & \hat{\mathcal{T}}^{+}_{g,\nu} \arrow[r] & \mathcal{N}_{g,\nu} \arrow[r] & \mathfrak{S}_{\nu} \\
\mathfrak{S}_{g,\nu} \arrow[r] \arrow[u] & \hat{\mathfrak{S}}_{g,\nu} \arrow[r] \arrow[u] & \mathcal{M}_{g,\nu} \arrow[u] &
\end{tikzcd}

\note{La flèche $\mathcal{N}_{g,\nu}\to\mathfrak{S}_{\nu}$ est encadrée~;
la lecture de son but est douteuse.}

d'où, en prenant les images, des groupes (8)
$\mathcal{T}^{+!}_{g,\nu}$, $\hat{\mathcal{T}}^{+!}_{g,\nu}$,
$\mathcal{N}^{!}_{g,\nu}$~; $\mathfrak{S}^{+!}_{g,\nu}$,
$\hat{\mathfrak{S}}^{+!}_{g,\nu}$, $\mathcal{M}^{!}_{g,\nu}$
s'interprétant comme des $\pi_1$ \ill{}…

Il y a lieu des variantes «~en~!~» (2')~(3')~(5'). \struck{TSVP} TSVP~!

\page{13}
(9)

\begin{tikzcd}[column sep=small, row sep=small, nodes={font=\scriptsize}]
1 \arrow[r] & \hat\pi_{g,\nu} \arrow[r] \arrow[d, no head, "\parallel"] & \mathcal{M}_{g,\nu} \arrow[r] \arrow[d, "?"'] & \mathcal{N}_{g,\nu} \arrow[r] \arrow[d, "?"'] & 1 \\
1 \arrow[r] & \hat\pi_{g,\nu} \arrow[r] & \mathrm{Aut}_{\mathrm{loc}}(\hat\pi_{g,\nu}) \arrow[r] & \mathrm{Autext}_{\mathrm{loc}}(\hat\pi_{g,\nu}) \arrow[r] & 1
\end{tikzcd}

prolongeant (2) et (3). On a donc que l'image $\mathcal{M}'_{g,\nu}$ de
$\mathcal{M}_{g,\nu}$ dans $\mathrm{Aut}_{\mathrm{loc}}(\hat\pi_{g,\nu})$
normalise celle $\hat{\mathcal{T}}'_{g,\nu}$ de
$\hat{\mathcal{T}}^{+}_{g,\nu}$ (le \uncertain{partie} «~Teichmüller
profini~») et l'image
\[
  (10)\qquad \mathcal{M}'_{g,\nu}/\mathcal{T}'_{g,\nu}
  \simeq \text{quotient de } \Gamma_{\bar{\mathbb{Q}}/\mathbb{Q}}
\]

\textbf{NB} on peut prouver que pour $g$ fixé, $\nu$ grand, c'est un
isom.\ avec $\Gamma_{\bar{\mathbb{Q}}/\mathbb{Q}}$ \uncertain{lui-même}
(\uncertain{\ill{}}). Mais je pense que c'est \uncertain{toujours} un
isom. C'est le cas pour $g=0$ (où le cas critique est $\nu=3$) -- ce qui
\struck{\ill{}} \add{\uncertain{équivaut à dire aussi} que}

\begin{tikzcd}
\Gamma_{\bar{\mathbb{Q}}/\mathbb{Q}} \arrow[r, hook] & \mathrm{Autext}_{\mathrm{loc}}(\hat\pi_{0,3})
\end{tikzcd}

\note{Sous cette flèche, une première version biffée~:
$\mathcal{M}_{0,3}\hookrightarrow\mathrm{Aut}$.}

(NB
\[
  \mathcal{N}_{0,3} \simeq \hat{\mathcal{T}}^{+}_{0,3} \times \mathcal{N}^{!}_{0,3},
  \qquad
  \hat{\mathcal{T}}^{+}_{0,3} = \mathcal{T}^{+}_{0,3} = \mathfrak{S}_3,
  \qquad
  \mathcal{N}^{!}_{0,3} \simeq \Gamma_{\bar{\mathbb{Q}}/\mathbb{Q}}\,)
\]
et pour $g=1$ (cas critique $\nu=1$)

\page{14}
\note{Il numérote ce feuillet (2) en haut à gauche~; la p.~12 en est
vraisemblablement le (1).}

Soit maintenant $k$ une ext.\ \emph{alg.\ close} de $\mathbb{Q}$
(\uncertain{\ill{}} $\bar{\mathbb{Q}}/\mathbb{Q}$), et $X$ une courbe de
type $g,\nu$ sur $k$, $x\in X$. Alors $X$ correspond à un pt géom.
\[
\begin{gathered}
  \xi\in\operatorname{Ob}\mathrm{Pt}(M_{g,\nu},k)
    \simeq\operatorname{Ob}\mathrm{Pt}(M_{g,\nu\,k}/k)\\
  X\simeq(\mathcal{X}_{g,\nu\,k})_\xi
\end{gathered}
\]
\struck{$\xi$,} $\{x\}\in\operatorname{Ob}\mathrm{Pt}(\mathcal{X}_{g,\nu\,k},k)$
${}\simeq\operatorname{Ob}\mathrm{Pt}(\mathcal{X}_{g,\nu\,k}/k)$
($x$ au-dessus de $\xi$)
\note{Dans les $\mathrm{Pt}$ de cette page comme de la p.~12, le symbole est
souligné sur la page.}

\struck{\ill{}} Considérons \struck{$\pi_1(\mathcal{X}_{g,\nu})$} (11)
\[
\begin{cases}
  \mathcal{T}_X \overset{\mathrm{déf}}{=} \pi_1(M_{g,\nu\,k},\xi),
    \quad \mathcal{M}_X=\pi_1(M_{g,\nu},\xi)\\
  \mathfrak{S}_{X,x} = \pi_1(\mathcal{X}_{g,\nu\,k},x),
    \quad \mathcal{N}_{X,x}=\pi_1(\mathcal{X}_{g,\nu},x)\\
  \pi_{X,x} = \pi_1(X,x)
\end{cases}
\]
\note{Les lettres $\mathcal{M}$ et $\mathcal{N}$ semblent ici échangées
par rapport à la p.~12 (où $\mathcal{N}$ est le $\pi_1$ de l'espace de
modules et $\mathcal{M}$ celui de la courbe universelle)~; le diagramme
(12) ci-dessous revient à l'usage de la p.~12. Transcrit tel quel.}

d'où (12)

\begin{tikzcd}[column sep=small, row sep=small, nodes={font=\scriptsize}]
 & & 1 \arrow[d] & 1 \arrow[d] & \\
1 \arrow[r] & \pi_{X,x} \arrow[r] \arrow[d, no head, "\parallel"] & \mathfrak{S}_{X,x} \arrow[r] \arrow[d] & \mathcal{T}_X \arrow[r] \arrow[d] & 1 \\
1 \arrow[r] & \pi_{X,x} \arrow[r] & \mathcal{M}_{X,x} \arrow[r] \arrow[d] & \mathcal{N}_X \arrow[r] \arrow[d] & 1 \\
 & & \Gamma_{\bar{\mathbb{Q}}/\mathbb{Q}} \arrow[r, "\sim"] & \Gamma_{\bar{\mathbb{Q}}/\mathbb{Q}} \arrow[d] & \\
 & & & 1 &
\end{tikzcd}

\marginal{\ill{} $\bar{\mathbb{Q}}$ \ill{} clôture alg.\ de $\mathbb{Q}$
dans $k$}

\note{Une première rédaction biffée commençait la deuxième ligne par
«~On a des isom.\ \ill{}~»~; au bas de la colonne de gauche, «~$\Gamma$~»
porte une marque biffée.}

On a un isom.\ \struck{\ill{}} extérieur \uncertain{canonique}
\struck{\ill{}}
\[
  (13)\qquad \mathcal{M}_{X,x}\xrightarrow{\ \sim\ }\mathcal{M}_{g,\nu}
\]
(i.e.\ isom.\ défini à automorph.\ int.\ près), définissant un isom.\ des
suites exactes entre (12) et (5),
\struck{conséquence~: (14)}

\note{Le diagramme suivant est biffé d'un trait sur chacune de ses lignes.}
\struck{$1\to\pi_{X,x}\to\mathfrak{S}_{X,x}\to\mathcal{T}_X\to 1$}

\struck{$1\to\hat\pi_{g,\nu}\to\hat{\mathfrak{S}}_{g,\nu}\to\hat{\mathcal{T}}_{g,\nu}\to 1$}
\note{Les deux lignes sont reliées par des flèches verticales marquées
$\simeq$.}
\[
  (14)\qquad \pi_{X,x}\ \overset{\mathrm{iso}}{\simeq}\ \hat\pi_{g,\nu}
\]
défini \emph{mod.\ opération} de
$\mathcal{M}_{g,\nu}\supset\hat{\mathfrak{S}}_{g,\nu}$.
C'est une structure supplémentaire \add{\uncertain{canonique}} sur le
groupe profini \uncertain{extérieur} $\pi_{X,x}$, \uncertain{Weil}
\struck{\ill{}} «~\uncertain{arithmétisation} \uncertain{canonique}~»~:
\uncertain{restriction} du groupe structural

\page{15}
$\mathrm{Autext}_{\mathrm{loc}}(\hat\pi_{g,\nu})$ dans
$\mathrm{Isomext}_{\mathrm{loc}}(\hat\pi_{g,\nu},\pi_{X,x})$, \ill{}
\uncertain{sous-groupe} $\mathcal{N}'_{g,\nu}$ \uncertain{image} de
$\mathcal{N}_{g,\nu}$.
\note{Sous $\mathrm{Autext}_{\mathrm{loc}}(\hat\pi_{g,\nu})$ la page porte
un signe $\cup$ au-dessus de $\mathcal{N}'_{g,\nu}$~; le début de la phrase
est sur la page précédente.}

Mais si $\mathcal{N}_{g,\nu}\xrightarrow{\mathrm{épi}}\mathcal{N}'_{g,\nu}$
n'était \uncertain{pas} iso, il serait quand${}={}$même plus fin de dire
que \uncertain{via} (13), on a un torseur \uncertain{canonique}
\[
  (15)\qquad
  P_{X,x} = \mathrm{Isom}_{\substack{\text{groupoïde}\\ \text{fond.\ de}\\ \mathcal{X}_{g,\nu\,\mathbb{Q}}}}(x_0,x)
\]
torseur à droite sous $\mathcal{M}_{g,\nu}$, et un hom.
\[
  (16)\qquad P_{X,x}\longrightarrow\mathrm{Isom}_{\mathrm{loc}}(\hat\pi_{g,\nu},\pi_{X,x})
\]
compatible avec
\[
  \struck{(17)}\qquad \mathcal{M}_{g,\nu}\xrightarrow{\ \mathrm{can}\ }\mathrm{Aut}_{\mathrm{loc}}(\hat\pi_{g,\nu})
\]
\note{Le but de cette flèche est récrit sur un premier mot biffé.}
sur les groupes d'opérateurs, d'où par passage au quotient par
$\hat\pi_{g,\nu}$
\[
  (17)\qquad Q_{X}\longrightarrow\mathrm{Isomext}_{\mathrm{loc}}(\hat\pi_{g,\nu},\pi_{X,x})
\]
$\Bigl(P_{X,x}/\hat\pi_{X,x}\simeq\mathrm{Isom}_{\substack{\text{groupoïde}\\ \text{fond.\ de } M_{g,\nu\,\mathbb{Q}}}}(\xi_0,\xi)\Bigr)$,
\note{Sous $Q_X$, un signe $\simeq$ vertical renvoie à cette parenthèse.}
compatible avec
\[
  \mathcal{N}_{g,\nu}\longrightarrow\mathrm{Autext}_{\mathrm{loc}}(\hat\pi_{g,\nu},\pi_{X,x})
\]
\struck{Et (17)} Les données (16) et (17) sont équivalentes. (Ce \ill{}
que structure suppl.\ sur $\pi_{X,x}$) -- c'est une «~arithmétisation~»
en un sens plus fin que \uncertain{via} $P'_{X,x}$, $Q'_{X}$ (images de
$P_{X,x}$, $Q_X$)

\page{16}
\note{Feuillet numéroté (3) de sa main, à la suite des p.~12 et~14.}

\textbf{NB} Si on \struck{réduit} \add{étend} le groupe structural de
$Q_X$ sous $\mathcal{N}_{g,\nu}$ via $\Gamma_{\bar{\mathbb{Q}}/\mathbb{Q}}$
(\add{\uncertain{sous-groupe} \ill{}}), on trouve
\[
  (18)\qquad R_X \overset{\mathrm{déf}}{=} Q_X \wedge^{\mathcal{N}_{g,\nu}} \Gamma_{\bar{\mathbb{Q}}/\mathbb{Q}}
  \simeq R_X/\hat{\mathcal{T}}^{+}_{g,\nu}
\]
\note{Le membre de droite se lit bien $R_X/\hat{\mathcal{T}}^{+}_{g,\nu}$,
sans doute pour $Q_X/\hat{\mathcal{T}}^{+}_{g,\nu}$~; transcrit tel quel.}
et ce $\Gamma_{\bar{\mathbb{Q}}/\mathbb{Q}}$-torseur est canonique~:
\[
  (19)\qquad
  \begin{cases}
    R_X \simeq \mathrm{Isom}_{\substack{\text{grpd.\ fond.}\\ \text{de } \operatorname{Spec}\mathbb{Q}}}(\operatorname{Spec}\bar{\mathbb{Q}},\operatorname{Spec}\tilde{\mathbb{Q}})\\
    \phantom{R_X}\simeq \mathrm{Isom}(\bar{\mathbb{Q}},\tilde{\mathbb{Q}})\\
    \phantom{R_X}\simeq \mathrm{plongts}(\bar{\mathbb{Q}},k)
  \end{cases}
\]
Donc, la donnée d'une trivialisation de $R_X$, i.e.\ d'une restriction
du groupe structural de $Q_X$ de $\mathcal{N}_{g,\nu}$ à
$\hat{\mathcal{T}}^{+}_{g,\nu}$ (ou de $P_{X,x}$ de $\mathcal{M}_{g,\nu}$
à $\hat{\mathfrak{S}}^{+}_{g,\nu}$) équivaut à la donnée d'un isom.
\[
  (20)\qquad \bar{\mathbb{Q}}\simeq\tilde{\mathbb{Q}}, \quad\text{i.e.\ plongement } \bar{\mathbb{Q}}\to k.
\]
\note{$\tilde{\mathbb{Q}}$ désigne vraisemblablement la clôture
algébrique de $\mathbb{Q}$ dans $k$, comme dans la marge de la p.~14.}

\page{17}
\textbf{NB} Donnons-nous un plongement
\[
  k\overset{\varphi}{\hookrightarrow}\mathbb{C}
\]
\struck{alors i.e.} \add{alors ce} plongement permet de réduire le groupe
structural du torseur $P_{X,x}$ \struck{\ill{}} \add{\uncertain{de}}
$\mathcal{M}_{g,\nu}$ au sous-groupe $\mathfrak{S}^{+}_{g,\nu}$, ce qui
revient au même, le groupe structural $\mathcal{N}_{g,\nu}$ de $Q_X$ au
sous-groupe $\mathcal{T}^{+}_{g,\nu}$ (structure de «~discrétification~»
\add{«~\uncertain{variantes}~»} des groupes profinis $\pi_{X,x}$)~:
\struck{a fortiori}
\[
  (21)\qquad
  \begin{cases}
    P^{\varphi}_{X,x}\subset P_{X,x}
      &\Longleftrightarrow\quad Q^{\varphi}_X\subset Q_X\\
    \mathfrak{S}^{+}_{g,\nu}\subset\mathcal{M}_{g,\nu}
      &\phantom{\Longleftrightarrow}\quad \mathcal{T}^{+}_{g,\nu}\subset\mathcal{N}_{g,\nu}
  \end{cases}
\]
a fortiori, par passage aux adhérences, on trouve
\[
  (22)\qquad
  \begin{cases}
    \widehat{P^{\varphi}_{X,x}}\subset P_{X,x}
      &\Longleftrightarrow\quad \hat{Q}^{\varphi}_X\subset Q_X\\
    \hat{\mathfrak{S}}^{+}_{g,\nu}\subset\mathcal{M}_{g,\nu}
      &\phantom{\Longleftrightarrow}\quad \hat{\mathcal{T}}^{+}_{g,\nu}\subset\mathcal{N}_{g,\nu}
  \end{cases}
\]
\struck{qui} \add{\uncertain{dont}} la donnée équivaut, par ce qui précède
à la page précédente, à celle d'un iso \struck{plongt}
$\tilde{\mathbb{Q}}\to\bar{\mathbb{Q}}$.
\note{Un double trait vertical, dans la marge gauche, marque les quatre
dernières lignes de la page.}
Cela montre que cette donnée (22) ne permet pas de reconstituer
\struck{\ill{}} \uncertain{la} donnée $\varphi$ elle-même, mais seulement
$\varphi\,|\,\tilde{\mathbb{Q}}$. Mais je pense que la connaissance de
\struck{$\varphi$} $P^{\varphi}_{X,x}$ (ou $Q^{\varphi}_X$) lui-même,
pour tout $X$ (et tout $g,\nu$) définit $\varphi$ ….

\page{18}
4) Récapitulation (II)
\note{Le feuillet porte «~4)~» en haut à gauche, à la suite des feuillets
numérotés (2) et (3), p.~14 et~16.}

$k$ corps \emph{alg.\ clos}, car.~0, i.e.\ \struck{\ill{}}
$\mathbb{Q}\to k$, $X$ courbe de type $g,\nu$ sur $k$, $x\in X$, d'où
\[
  \xi\in M_{g,\nu}(k)=\Gamma(M_{g,\nu\,k}/k)
  \qquad X=\mathcal{X}_{g,\nu\,\xi}=(\mathcal{X}_{g,\nu\,k})_\xi,
  \qquad x\in\mathcal{X}_{g,\nu}(k)
\]
$\xi$, $x$ pts géométriques de $M_{g,\nu\,\mathbb{Q}}$,
$\mathcal{X}_{g,\nu\,\mathbb{Q}}$, d'où par fibration exacte
\note{Dans la formule ci-dessus, le $\Gamma$ est souligné sur la page~; un
signe biffé, illisible, précède $M_{g,\nu}(k)$.}

(0)

\begin{tikzcd}[column sep=small, row sep=small, nodes={font=\scriptsize}]
\mathcal{X}_{g,\nu\,\mathbb{Q}} \arrow[d] & \mathcal{X}_{g,\nu\,k} \arrow[l, no head] \arrow[d, no head] \\
M_{g,\nu\,\mathbb{Q}} \arrow[d, no head] & M_{g,\nu\,k} \arrow[d, no head] \\
\mathbb{Q} & k \arrow[l, no head]
\end{tikzcd}

le diagramme de suites exactes

(1)

\begin{tikzcd}[column sep=small, row sep=small, nodes={font=\scriptsize}]
 & 1 \arrow[d] & 1 \arrow[d] & & \\
 & \pi_{X,x} \arrow[r, no head, "="] \arrow[d] & \pi_{X,x} \arrow[d] & & \\
1 \arrow[r] & \mathfrak{S}_{X,x} \arrow[r] \arrow[d] & M_{X,x} \arrow[r] \arrow[d] & \Gamma_{\tilde{\mathbb{Q}}/\mathbb{Q}} \arrow[r] \arrow[d, no head, "\parallel"] & 1 \\
1 \arrow[r] & \mathcal{T}_X \arrow[r] \arrow[d] & N_X \arrow[r] \arrow[d] & \Gamma_{\tilde{\mathbb{Q}}/\mathbb{Q}} \arrow[r] & 1 \\
 & 1 & 1 & &
\end{tikzcd}

\marginal{$\tilde{\mathbb{Q}}\to\mathbb{C}$ clôture alg.\ de
$\mathbb{Q}$ dans $k$}

\textbf{NB} Le diagramme équivaut à la donnée
\struck{\ill{}} \add{de $N_X$}
\[
  1\subset\pi_{X,x}\subset\mathfrak{S}_{X,x}\subset M_{X,x},
\]
\note{Sous la chaîne, deux accolades~: l'une sous
$\mathfrak{S}_{X,x}/\pi_{X,x}$ marquée $\mathcal{T}_X$, l'autre sous
$M_{X,x}/\mathfrak{S}_{X,x}$ marquée $\Gamma_{\tilde{\mathbb{Q}}/\mathbb{Q}}$.}
où
\[
  (2)\qquad \text{déf}\quad
  \begin{cases}
    \pi_{X,x}=\pi_1(X,x) & (\simeq\hat\pi_{g,\nu}\ \text{\uncertain{non can.}})\\
    \mathfrak{S}_{X,x}=\pi_1(\mathcal{X}_{g,\nu\,k},x) & (\simeq\hat{\mathfrak{S}}_{g,\nu}\ \text{---})\\
    \mathcal{T}_X=\pi_1(M_{g,\nu\,k},\xi) & (\simeq\hat{\mathcal{T}}_{g,\nu}\ \cdots)\\
    M_{X,x}=\pi_1(\mathcal{X}_{g,\nu},x) & (\simeq M_{g,\nu}\ \cdots)\\
    N_X=\pi_1(M_{g,\nu},\xi) & (\simeq N_{g,\nu}\ \cdots)
  \end{cases}
\]
\note{Sur cette page les groupes $M$, $N$ sont écrits en capitales
droites, non en ronde comme $\mathcal{M}$, $\mathcal{N}$ aux p.~12--17~;
transcrits tels quels. Devant (2), une première amorce biffée
(«~$\mathfrak{S}_{X,x}={}$~»).}

Les derniers membres $\hat\pi_{g,\nu}$ \uncertain{etc.} se définissent
comme les premiers, mais à l'aide de $(X_0,x_0)$ sur $\mathbb{C}$, de type
$g,\nu$, \uncertain{définissant} les pts géom.\ $x_0$ et $\xi_0$ de
\struck{\ill{}} $\mathcal{X}_{g,\nu\,\mathbb{Q}}$ et $M_{g,\nu\,\mathbb{Q}}$.

\page{19}
Les \uncertain{choix} des derniers \uncertain{pts} $\xi_0$~:
\uncertain{$\xi$} définissent un torseur
\[
  (3)\qquad P_{X,x}\ \text{torseur à droite sous } \mathcal{M}_{g,\nu}
  \quad(\text{\uncertain{commutant} à } \mathcal{M}_{X,x})
\]
de telle façon que (1) soit \uncertain{déduit} du diagramme analogue pour
$(X_0,x_0)$

($1_{g,\nu}$)

\begin{tikzcd}[column sep=small, row sep=small, nodes={font=\scriptsize}]
 & 1 \arrow[d] & 1 \arrow[d] & & \\
 & \hat\pi_{g,\nu} \arrow[r, no head, "="] \arrow[d] & \hat\pi_{g,\nu} \arrow[d] & & \\
1 \arrow[r] & \hat{\mathfrak{S}}_{g,\nu} \arrow[r] \arrow[d] & \mathcal{M}_{g,\nu} \arrow[r] \arrow[d] & \Gamma_{\bar{\mathbb{Q}}/\mathbb{Q}} \arrow[r] \arrow[d, no head, "\parallel"] & 1 \\
1 \arrow[r] & \hat{\mathcal{T}}_{g,\nu} \arrow[r] \arrow[d] & \mathcal{N}_{g,\nu} \arrow[r] \arrow[d] & \Gamma_{\bar{\mathbb{Q}}/\mathbb{Q}} \arrow[r] & 1 \\
 & 1 & 1 & &
\end{tikzcd}

par le torseur, grâce à l'action de $\mathcal{M}_{g,\nu}$ opérant
\uncertain{par automorphismes sur lui-même}, et sur sa
\uncertain{filtration} par $\hat{\mathfrak{S}}_{g,\nu}$, $\hat\pi_{g,\nu}$.
En particulier
\[
  (4)\qquad
  \begin{cases}
    \pi_{X,x}\simeq P_{X,x}\wedge^{\mathcal{M}_{g,\nu}}\hat\pi_{g,\nu}\\
    \mathfrak{S}_{X,x}\simeq P_{X,x}\wedge^{\mathcal{M}_{g,\nu}}\hat{\mathfrak{S}}_{g,\nu}\\
    \mathcal{T}_{X,x}\simeq P_{X,x}\wedge^{\mathcal{M}_{g,\nu}}\hat{\mathcal{T}}_{g,\nu}
      \simeq Q_X\wedge^{\mathcal{N}_{g,\nu}}\hat{\mathcal{T}}_{g,\nu}\\
    \mathcal{M}_{X,x}\simeq P_{X,x}\wedge^{\mathcal{M}_{g,\nu}}\mathcal{M}_{g,\nu}\\
    \mathcal{N}_X\simeq P_{X,x}\wedge^{\mathcal{M}_{g,\nu}}\mathcal{N}_{g,\nu}
      \simeq Q_X\wedge^{\mathcal{N}_{g,\nu}}\mathcal{N}_{g,\nu}\\
    \tilde{\mathbb{Q}}\simeq P_{X,x}\wedge^{\mathcal{M}_{g,\nu}}\bar{\mathbb{Q}}
      \simeq Q_X\wedge^{\mathcal{N}_{g,\nu}}\bar{\mathbb{Q}}
      \simeq R\wedge^{\Gamma_{\bar{\mathbb{Q}}/\mathbb{Q}}}\bar{\mathbb{Q}}
  \end{cases}
\]
\note{Les produits contractés sont écrits $P\wedge_{G}A$, le groupe en
indice~; on les rend $P\wedge^{G}A$. À la fin de la deuxième ligne, un
$\mathcal{M}_{X,x}$ biffé.}

\page{20}
\note{Feuillet marqué «~5)~» en haut à gauche.}

où
\[
  (5)\qquad Q_X=P_{X,x}\wedge^{\mathcal{M}_{g,\nu}}\mathcal{N}_{g,\nu},
  \qquad
  R=P_{X,x}\wedge^{\mathcal{M}_{g,\nu}}\Gamma_{\bar{\mathbb{Q}}/\mathbb{Q}}
  \simeq Q_X\wedge^{\mathcal{N}_{g,\nu}}\Gamma_{\bar{\mathbb{Q}}/\mathbb{Q}}
\]
Ainsi, $\tilde{\mathbb{Q}}$ et (1) sont connus par la connaissance du
\add{$\mathcal{M}_{g,\nu}$-tors.} $P_{X,x}$ (3). Mais bien sûr, même si
$k$ est alg.\ sur $\mathbb{Q}$ ($k=\tilde{\mathbb{Q}}$, voire
$k=\bar{\mathbb{Q}}$ i.e.\ avec $R$ trivialisé) la donnée de
\struck{$Q$} $P_{X,x}$ ne permet pas de \uncertain{récupérer} $X$ sur $k$
-- p.\,ex.\ les automorph.\ des torseurs $P_{X,x}$ sont les éléments de
$\mathcal{M}_{X,x}$ -- \uncertain{ils} \uncertain{ont la} puissance du
continu -- quand \uncertain{même} \struck{$\tilde{\mathbb{Q}}$}
$\hat{\mathfrak{S}}_{g,\nu}$-torseurs correspondant à une restriction des
groupes structuraux à $\hat{\mathcal{T}}_{g,\nu}$ i.e.\ à un iso.\
$\tilde{\mathbb{Q}}\simeq\bar{\mathbb{Q}}$, correspondant à un élément de
$\mathfrak{S}_{X,x}$ -- il a la puissance du continu, tandis que les
\uncertain{automorph.}\ de $X,x$ compatibles avec l'identité sur
$k=\tilde{\mathbb{Q}}$ (${}\simeq\bar{\mathbb{Q}}$ \add{donné}) forment
un groupe \emph{fini}~!

Mais …
\[
  (6)\qquad k=\varinjlim_i K_i
\]
les $K_i$ sous-corps de t.f.\ -- et \struck{\ill{}} \add{la} structure
\uncertain{dite} $(X,x)$ implique une structure
\note{L'argument se poursuit au-delà de ce lot.}

\end{document}
